\documentclass{resume}

\usepackage[left=0.4in,top=0.4in,right=0.4in,bottom=0.4in]{geometry}
\usepackage{fontawesome5}
\newcommand{\tab}[1]{\hspace{.2667\textwidth}\rlap{#1}}
\newcommand{\itab}[1]{\hspace{0em}\rlap{#1}}
\name{Kevin Lee}
\address{+1 530-376-9732 \\ Davis, CA}
\address{\small\faEnvelope\ \href{mailto:yhylee@ucdavis.edu}{yhylee@ucdavis.edu} \quad
\faLinkedin\ \href{https://www.linkedin.com/in/kevin-lee-b78448388/}{linkedin.com/in/kevin-lee-b78448388} \quad
\faGithub\ \href{https://github.com/section9-us}{github.com/section9-us}}

\begin{document}

\begin{rSection}{PROFESSIONAL SUMMARY}
Security-focused Software Engineer with 10+ years of experience protecting public-sector and critical-infrastructure systems. Authored and disseminated 500+ cyber threat intelligence (CTI) reports annually; delivered penetration testing, red teaming, vulnerability analysis, and remediation consulting for 200+ public institutions. Combines software development and incident response with current UC Davis M.S. research applying AI and machine learning to cyber threat detection.
\end{rSection}

\begin{rSection}{EXPERIENCE}

\textbf{Security Engineer, Threat Management Division} \hfill Jan 2022 -- Present\\
Ministry of Science and ICT \hfill \textit{Sejong, Republic of Korea}
\begin{itemize}
    \itemsep -4pt {}
    \item Coordinated critical-infrastructure attack analysis; authored and disseminated 500+ CTI reports annually on malware, attack evidence, and APT activity, providing timely memoranda and remediation guidance to allied U.S., U.K., and European government agencies.
    \item Developed security patches and authored vulnerability advisories in approximately 50 cases over three years; disseminated formal remediation notices to 7,000+ South Korean government agencies and public institutions.
    \item Initiated a 2023 research collaboration with the Korea Internet \& Security Agency (KISA), analyzing international threat-group infrastructure and exploring AI and data analytics for critical-infrastructure attack prediction and mitigation.
\end{itemize}

\textbf{Security Engineer, Security Planning Division} \hfill May 2018 -- Jan 2022\\
Ministry of Science and ICT \hfill \textit{Sejong, Republic of Korea}
\begin{itemize}
    \itemsep -4pt {}
    \item Delivered on-site red teaming and security consulting across nuclear, thermal-power, railway-control, and power systems at 50+ public institutions annually (200+ total).
    \item Led on-site cyber incident-investigation teams of engineers seconded from the military, police, and research institutes, directing threat detection, initial response, and protective security measures.
    \item Partnered with KISA and Ministry teams to analyze several discovered vulnerabilities, map them to CWE categories, and coordinate CVE ID assignments.
\end{itemize}

\textbf{Software Engineer, Reactor Design Development Division} \hfill Nov 2014 -- Mar 2018\\
KEPCO E\&C \hfill \textit{Gimcheon, Republic of Korea}
\begin{itemize}
    \itemsep -4pt {}
    \item Developed reactor design software and an engineering database for nuclear-plant design workflows over 3+ years.
    \item Identified software vulnerabilities and performed secure coding remediation for reactor software and engineering databases in a high-reliability critical-infrastructure environment.
\end{itemize}

\end{rSection}

\begin{rSection}{EDUCATION}

{\bf Master of Science in Computer Science}, University of California, Davis \hfill {In Progress; Expected Mar 2027}\\
Current GPA: 4.0/4.0\\
Korean Government Long-term Fellowship for Overseas Study, Ministry of Science and ICT

{\bf Bachelor of Science in Computer Software}, Kwangwoon University (GPA: 3.32/4.0) \hfill {Feb 2015}

\end{rSection}

\newpage

\begin{rSection}{RESEARCH}

\textit{UC Davis M.S. research portfolio: AI-based threat detection, exposure-aware vulnerability triage, and critical-service resilience.}

\href{https://github.com/section9-us/Network_ids-neural-vs-traditional-attack-eval}{\textbf{Attack-Type-Specific Network Intrusion Detection}} \hfill Jun 2026
\begin{itemize}
    \itemsep -4pt {}
    \item Investigated whether aggregate model scores conceal attack-specific detection gaps; built a reproducible Python pipeline comparing traditional and PyTorch models on 50,000 sampled CSE-CIC-IDS2018 flows.
    \item Measured false negatives by attack type and family; the best model detected only 29.9\% of Infiltration attacks in the held-out sample. Documented sampling, class imbalance, and benchmark dependence rather than claiming production IDS performance.
\end{itemize}

\href{https://github.com/section9-us/Security_CVE-CWE-Network-Prioritizer}{\textbf{CVE/CWE Network Vulnerability Prioritizer}} \hfill May 2026
\begin{itemize}
    \itemsep -4pt {}
    \item Investigated how exposed-service context changes CVE triage; built a Python application for authorized scanning, fingerprinting, and NVD CVE/CWE correlation, ranking reachable services by CVSS severity, known-exploitation signals, and fingerprint confidence.
    \item Added a browser dashboard, report export, and SQLite scan history; packaged the tool with Docker and a GitHub Actions test workflow, with 29 passing tests in the documented local unit suite.
\end{itemize}

\href{https://github.com/section9-us/Security_Resilient-DNS-Against-DDoS}{\textbf{DNS Availability Under DoS}} \hfill Jan 2026
\begin{itemize}
    \itemsep -4pt {}
    \item Evaluated DNS resilience under synthetic DoS load and node failure with a Docker/Knot DNS testbed comparing a single server, round-robin distribution, and health- and latency-aware routing.
    \item Used Python orchestration and InfluxDB to measure query availability, latency, and error rates across 30 probe queries per deployment phase; documented the single-host scope rather than representing it as production BGP anycast.
\end{itemize}

\end{rSection}

\begin{rSection}{SKILLS}
\begingroup
\renewcommand{\arraystretch}{1.5}
\begin{tabular}{ @{} >{\bfseries}l @{\hspace{2ex}} p{0.77\textwidth} }
Programming & Python, Java, C/C++, C\#, JavaScript, HTML/CSS, SQL, Spring MVC \\
Security & Malware Analysis, CTI, APT Analysis, Incident Response, Red Teaming, Penetration Testing, Threat Modeling, Vulnerability Assessments, and Secure Coding \\
AI & Machine Learning, LLMs, Generative AI, Prompt Engineering, RAG, Model Evaluation, Security Automation, and MLOps \\
Tools & Burp Suite, Nmap, Wireshark, IDA Pro, PyTorch, TensorFlow, Scikit-learn, Hugging Face, MySQL, PostgreSQL, Oracle, and Docker \\
Languages & Korean (Native), English (Fluent) \\
\end{tabular}\\
\endgroup
\end{rSection}

\begin{rSection}{CERTIFICATIONS \& TRAINING}
\textbf{Certifications:} Engineer Information Processing, HRDK (2013) \quad SQL Developer, Korea Data Agency (2014) \\[0.5\baselineskip]
\textbf{Training:} Locked Shields, NATO CCDCOE (2023) \quad Advanced Web Application Exploits, Black Hat USA (2021)
\end{rSection}

\end{document}
